\section{Threat Model and Problem Formulation}
\label{sec:model}

\subsection{Entities, Trust and Adversary}

A market instance comprises subjects $\mathcal{S}$ with $|\mathcal{S}|=N$,
platforms $\mathcal{P}$ with $|\mathcal{P}|=n_p$, tenants $\mathcal{T}$, and an
append-only ledger $\mathcal{L}$. Each subject $s$ contributes records to one or more platforms and carries a
declared cap $E_s$ on cumulative privacy loss, with an SLO saying how the cap may
be spent over time. Each tenant submits queries naming a predicate and the
platforms to evaluate it on; the receiving platform prices the query, obtains an
admission decision, evaluates the noisy answer, and settles the charge.

The coordinator is \emph{crash-faulty but not Byzantine}: it may stop and lose
volatile state, but it does not fabricate charges. The ledger is append-only and
authenticated; anchoring it on a public chain strengthens non-repudiation without
changing any argument here. Platform code is assumed to
compute the declared mechanism, with attestation left to orthogonal
infrastructure~\cite{liu2025yotta,zhu2025auditing}. $\mathrm{supp}(q)$ follows from the predicate and the platform's
record index, never from private values, so a refusal reveals only that a budget is
spent. Out of scope are platform--subject collusion, timing side channels, and a
platform that answers with less noise than it charges for, which is a
verifiable-computation problem.

The adversary is a tenant controlling $k\ge1$ accounts, possibly over several
platforms, choosing queries adaptively as a function of the answers received, and
aiming to extract more about a target subject than its cap allows. Because the
ledger charges accounts rather than subjects, its natural strategy is
\emph{splitting}: divide a sequence that together exceeds $E_s$ into $k$ fragments,
each within the per-account allowance, and issue them under distinct identities.
Splitting is what a rational tenant does when the price of a query is sub-linear in
the information it yields. The adversary may also crash the coordinator: charging
before execution and crashing loses the charge, charging after execution and
crashing delivers the answer for free. Either direction breaks the accounting
without any cryptographic failure, which is why settlement is part of the privacy
mechanism rather than a billing detail.

\subsection{Cost Model}

Every query $q$ carries a cost vector
$\cost(q)=(\eps_{\mathrm{pure}},\rho,\delta)$ describing the loss charged to
\emph{each} subject in its support. For a Gaussian mechanism with sensitivity
$\Delta$ and scale $\sigma$, $\rho=\Delta^{2}/(2\sigma^{2})$ and
$\eps_{\mathrm{pure}}=0$; for a Laplace mechanism with scale $b$,
$\eps_{\mathrm{pure}}=\Delta/b$ and $\rho=0$. Charges add in $\rho$ under zCDP
composition and in $\eps_{\mathrm{pure}}$ under pure DP composition, so the cost of
a transcript is the vector sum of its queries. Write $\spend_s(t)$ for the
cumulative vector charged to $s$ and
\begin{equation}
\spend_s(t,W)=\sum_{q\,:\,s\in\mathrm{supp}(q),\; t_q\in(t-W,\,t]} \cost(q)
\label{eq:window}
\end{equation}
for the spend inside a window of length $W$. The hard cap is scalar: a cumulative
vector $c$ is within budget when $g(c)\le E_s$, where
$g(c)=\eps_{\mathrm{pure}}+\rho+2\sqrt{\rho\ln(1/\delta)}$ converts a pure-DP
part and a zCDP part to $\eps$ at the market's $\delta$. Both mechanisms we price
carry $\delta{=}0$ in the vector---Laplace is pure, Gaussian is zCDP---so the
conversion consumes $\delta$ once rather than once per query. The test must be on
$g$, not on the two components: at a $1.0\eps$ cap, admitting while
$\eps_{\mathrm{pure}}\le1$ and $\rho\le0.021$ separately lets a mixed stream reach
a converted $1.41\eps$, whereas the scalar test rejects the 13 queries that cause it
(E24). The SLO constrains the
windowed spend under the same scalarisation.

\subsection{The Enforcement Problem}

An enforcement mechanism is a function
$a:(q,t)\mapsto\{\textsc{admit},\textsc{throttle},\textsc{reject}\}$ evaluated
before $q$ executes. Four properties are required.

\emph{Soundness.} Write $\cost_{\mathrm{exec}}(q)$ for the cost of the mechanism
actually executed: $\cost(q)$ in \textsc{normal} mode and $\beta\cost(q)$ in
\textsc{throttle} mode, where $\beta\in(0,1]$ is the charge scale of
Section~\ref{sec:design-slo}. If $a(q,t)\in\{\textsc{admit},\textsc{throttle}\}$
then for every $s\in\mathrm{supp}(q)$,
$g(\spend_s(t)+\cost_{\mathrm{exec}}(q))\le E_s$, and no executed transcript ever
produces $g(\spend_s(t))>E_s$. The two modes are covered by one
condition because a throttled query is still answered: what changes is the noise,
and the charge follows the noise. Soundness is what
distinguishes enforcement from auditing: it is a statement about the transcript
that is never released, not about a sample of past releases.

\emph{Availability.} Let $A$ be the admitted fraction and $U$ the total value of
admitted answers. A mechanism should maximise $A$ and $U$ subject to soundness,
and in particular should not refuse a query that is provably safe.

\emph{Stability.} Let $\tau$ be the number of mode transitions per epoch.
Oscillation between admitting and refusing is a cost not captured by $A$, because
a tenant that cannot predict admission cannot plan a workload; we require $\tau$
to be bounded as a function of the hysteresis and dwell parameters.

\emph{Atomicity.} For every execution, including those with crashes and
concurrent submissions, the charge committed to a subject equals the sum of costs
of the queries admitted for that subject, each query is charged at most once, and
no in-flight reservation is counted as available capacity.

The SLO specification is $(\theta_{\mathrm{soft}},p,W)$ with
$\theta_{\mathrm{soft}}=\alpha E_s$ for a soft ratio $\alpha\in(0,1)$: the
fraction of subject-epochs whose windowed spend exceeds $\theta_{\mathrm{soft}}$
must not exceed $p$, and no subject-epoch may exceed $E_s$. The two thresholds
are deliberately separate. The soft threshold is a target the breaker may miss
under adversarial load and $p$ bounds how often; the hard cap is never a target,
and one violation is a failure of the mechanism rather than a degradation of
service.

\subsection{Success Criterion for an Attacker}

To compare mechanisms under splitting we report the \emph{escape ratio}
\begin{equation}
R=\max_{s\in\mathcal{S}}\;\spend_s(T)/E_s ,
\label{eq:escape}
\end{equation}
where $T$ is the trace length and $\spend_s(T)$ the total charged to $s$ over
it: the worst-case ratio of loss to declared cap. $R\le1$ certifies soundness on the trace; $R>1$ quantifies the
leak. The ratio folds the effect of splitting, of per-account keying and of the
composition rule into one quantity comparable across mechanisms.
